\section{Normalization and domain audit}\label{app:normalization}

This appendix expands the interfaces that are especially vulnerable to a
factor, sign, or domain error.  It is part of the proof, rather than a report
about the accompanying computations.  The technical supplement records the
machine implementations of the same interfaces.

\subsection{The entire completed function}

Let $\Lambda_0$ denote the pole-removed entire completion used in the formal
library.  Its relation to the meromorphic completion is
\[
 \Lambda(s)=\Lambda_0(s)-\frac1s-\frac1{1-s}.
\]
The entire function in this paper is therefore
\begin{equation}\label{app:xi-entire}
 \xi(s)=\frac{s(s-1)}2\Lambda_0(s)+\frac12.
\end{equation}
For $s\ne0,1$, substitution of the preceding identity in
\eqref{app:xi-entire} cancels the two polar terms and gives the usual
$s(s-1)\Lambda(s)/2$.  Both sides are entire, so the equality extends across
$0$ and $1$.  The functional equation $\Lambda_0(1-s)=\Lambda_0(s)$ then
gives
\[
 \xi(1-s)=\xi(s).
\]
Consequently $w\mapsto\xi(1/2+w)$ is even and its odd derivatives at zero
vanish.

The coefficient convention is worth checking at the derivative level.  From
\[
 \xi(1/2+w)=\sum_{n\ge0}\frac{\gamma(n)}{n!}w^{2n}
\]
one obtains exactly
\begin{equation}\label{app:gamma-derivative}
 \gamma(n)=\frac{n!}{(2n)!}
 \left.\frac{\dd^{,2n}}{\dd w^{2n}}\xi(1/2+w)\right|_{w=0}.
\end{equation}
There is no factor $2^{2n}$ in \eqref{app:gamma-derivative}.  This explicit
derivative check is also an equation-regression gate in the release package.

\subsection{The factor-eight calculation}

Riemann's theta-kernel identity, after the substitution $t=e^{2u}$ and
symmetrization about the critical line, gives
\begin{equation}\label{app:factor-eight-derivative}
 \left.\frac{\dd^{,2n}}{\dd w^{2n}}\xi(1/2+w)\right|_{w=0}
 =8\int_0^\infty\omega(e^{2u})e^{u/2}u^{2n}\,\dd u.
\end{equation}
Combining \eqref{app:gamma-derivative} and
\eqref{app:factor-eight-derivative} gives the moment formula used in
Section~\ref{sec:mellin}.  The factor $e^{u/2}$ is the Jacobian and centered
power contribution after $t=e^{2u}$; replacing it by $e^u$ would change even
the zeroth moment.  The release regression evaluates the two sides
independently at low orders precisely to detect either transcription error.

For the later saddle calculation it is convenient to integrate by parts and
write
\[
 F(s)=\int_0^\infty\mathcal K(t)t^s\,\dd t.
\]
The boundary terms vanish because the theta kernel decays exponentially at
infinity and its modular transform supplies the corresponding decay at
zero.  Two integrations give
\begin{equation}\label{app:moment-assembly}
 \gamma(M)=\frac{\Gamma(M+1)}{\Gamma(2M+1)}2^{-2M-2}
 \left(32\binom{2M}{2}F(2M-2)-F(2M)\right).
\end{equation}
The two Mellin arguments in \eqref{app:moment-assembly} explain the shift
$N=2M-2$ used throughout the saddle analysis.
Define
\[
 M_z=2^{-2z-2}\left(32\binom{2z}{2}F(2z-2)-F(2z)\right),
\]
so that $\gamma(z)=\Gamma(z+1)M_z/\Gamma(2z+1)$ exactly.

\subsection{Nested sectors and proportional discs}

Fix
\[
 \mathcal S_1(R)=\{M:|M|>R,\ |\arg M|<1/200\},\qquad
 \overline{\mathcal S}_0(R)
 =\{M:|M|\ge R,\ |\arg M|\le1/400\}.
\]
All logarithms and square roots are continued from the positive axis in
$\mathcal S_1(R)$.  There is an $\eta_0>0$, depending only on the difference
between the two displayed angles, such that
\begin{equation}\label{app:sector-disc}
 M\in\overline{\mathcal S}_0(R),\quad |z-M|\le\eta_0|M|
 \quad\Longrightarrow\quad z\in\mathcal S_1(R/2).
\end{equation}
Indeed, $|z|\ge(1-\eta_0)|M|$, while the angular displacement is at most
$\arcsin(\eta_0/(1-\eta_0))$.  Choose $\eta_0$ so this is below $1/400$.
This elementary separation is the geometric input that permits Cauchy
differentiation of the relative error.  It also explains why the main-term
estimate may be stated on a wider sector than the final uniform error.

\section{Detailed saddle estimates}\label{app:saddle-details}

\subsection{Legality of the logarithmic equation}

Put $\ell=\log|N|$ and
\[
 L_0=\log N-\log\log N-\log\pi.
\]
On the closed disc $D_N=\{L:|L-L_0|\le1\}$ define
\begin{align*}
 m(\ell)&=\ell-\log(\ell+1)-\log\pi-1,\\
 M_\theta(\ell)&=\ell+2+\log(\ell+1)+\theta/\ell+\log\pi.
\end{align*}
For $\ell\ge12$ and $|\arg N|\le\theta$ the triangle inequality gives
\begin{equation}\label{app:disc-geometry}
 \Re L\ge m(\ell)>0,\qquad |L|\le M_\theta(\ell).
\end{equation}
Thus the principal logarithm of $L$ is legal on the whole disc.  Moreover,
\begin{equation}\label{app:saddle-small-term}
 \left|\frac{3L}{4N}\right|
 \le\frac{3M_\theta(\ell)}{4e^\ell}<10^{-4},
\end{equation}
so the correction logarithm is legal there as well.

For $|z|\le1/2$, integration along the segment from $0$ to $z$ gives
\[
 \log(1+z)=\int_0^1\frac{z}{1+tz}\,\dd t,
 \qquad |\log(1+z)|\le2|z|.
\]
Apply this first to $L/L_0-1$ and then to $-3L/(4N)$.  Since
$|L_0|\ge\ell/2$, the comparison function in the proof of
Lemma~\ref{lem:saddle} satisfies
\begin{align}
 \sup_{\partial D_N}|G_N(L)-(L-L_0)|
 &\le
 2\frac{\log(\ell+1)+\theta/\ell+\log\pi+1}{\ell}
 +\frac{3M_\theta(\ell)}{2e^\ell}.\label{app:rouche-bound}
\end{align}
The right side is decreasing from $\ell=12$ and is strictly below one at
that endpoint.  Rouch\'e therefore gives one zero counted with multiplicity.
No numerical sample of an interior zero is used in this conclusion.

\subsection{Simplicity before differentiation}

At the zero, exponentiation of the logarithmic equation yields
$N=L_N(\pi e^{L_N}+3/4)$.  Set $r=L_N^{-1}$ and $\sigma=L_N/N$.  The
derivative denominator factors as
\begin{equation}\label{app:Q-factor}
 Q_N=(1+L_N)N-\frac34L_N^2
 =NL_N\left(1+r-\frac34\sigma\right).
\end{equation}
After enlarging $R$, the disc geometry gives $|r|,|\sigma|\le7/50$, whence
\begin{equation}\label{app:Q-margins}
 \left|1+r-\frac34\sigma\right|
 \ge1-\frac7{50}-\frac{21}{200}=\frac{151}{200}>\frac9{16}.
\end{equation}
Thus the root is simple before the implicit-function theorem is invoked.
Local branches patch because any two branches meeting the same Rouch\'e disc
must select its unique zero.  Differentiation now legitimately gives
$L_N'=L_N/Q_N$.

The denominator arising after six differentiations is a power of
$4+4r-3\sigma$.  The same box gives the stronger margin
\begin{equation}\label{app:H6-margin}
 |4+4r-3\sigma|\ge4-\frac{28}{50}-\frac{21}{50}
 =\frac{151}{50}.
\end{equation}
Equations \eqref{app:Q-margins} and \eqref{app:H6-margin} are kept separate:
the first licenses analytic differentiation; the second controls the exact
rational derivative tower.

\subsection{Imaginary part and Gaussian direction}

Write $L_N=a+ib$ and $N=|N|e^{i\phi}$.  Taking arguments in the saddle
equation gives
\[
 \phi=\arg L_N+\arg(\pi e^{L_N}+3/4).
\]
Because $a\ge m(12)>7.29$, the second argument differs from $b$ by at most
$3e^{-a}/(2\pi)$, and $|\arg L_N|\le |b|/a$.  Therefore
\begin{equation}\label{app:imag-bootstrap}
 |b|\le|\phi|+\frac{|b|}{a}+\frac{3e^{-a}}{2\pi}.
\end{equation}
Solving \eqref{app:imag-bootstrap} at $|\phi|\le1/100$ gives
$|b|<0.012<1/20$.  This direct argument avoids deriving the imaginary bound
from the much coarser estimate $|L_N-L_0|<1$.

Finally,
\[
 K_N=\frac{N}{L_N}
 \left(1+\frac1{L_N}-\frac{3L_N}{4N}\right).
\]
The last factor is in the disc of radius $49/200$ about $1$.  Combining its
argument with \eqref{app:imag-bootstrap}, and increasing $R$ once more,
gives
\begin{equation}\label{app:K-sector}
 |\arg K_N|<1/20,\qquad
 \Re K_N\ge\cos(1/20)|K_N|.
\end{equation}
The right-half-plane inequality, not merely $K_N\ne0$, supplies the uniform
real Gaussian majorant.

\section{Contour and coefficient details}\label{app:contour-details}

\subsection{Truncation and connectors}

Set
\[
 \Phi_N(u)=N\operatorname{Log}u-\frac34u-\pi e^u,
 \qquad g_N(u)=e^u e^{\Phi_N(u)},
 \qquad I_1(N):=F_1(N):=\int_0^\infty g_N(u)\,\dd u.
\]
Write $L_N=\alpha+ib$ and first split this integral at $1$.
Apply Cauchy's theorem to the finite rectangle with vertices
$1,X,X+ib,1+ib$.  The right connector tends to zero as $X\to\infty$, so
\begin{equation}\label{app:legal-contour}
 I_1(N)=E_N+\int_1^\infty g_N(x+ib)\,\dd x,
 \qquad
 E_N=\int_0^1g_N(x)\,\dd x+i\int_0^b g_N(1+iy)\,\dd y.
\end{equation}
The horizontal ray lies in $\Re u\ge1$, hence in a fixed principal-log
domain.  The finite endpoint term satisfies
\begin{equation}\label{app:connector-small}
 |E_N|\le|\mathcal A(N)|e^{-c'|N|\log\log|N|}.
\end{equation}
This order of limits prevents the invalid shortcut of translating the whole
infinite real line across the logarithmic cut.

\subsection{Central window and odd cancellation}

On the actual shifted ray write $u=L_N+r$, so $r\ge1-\Re L_N$, and choose
$V=|K_N|^{-2/5}$.  For sufficiently large $|N|$, the central window is
contained in this ray.  Taylor's formula with integral remainder yields,
uniformly for $|r|\le V$,
\begin{equation}\label{app:phase-four}
 \Phi_N(L_N+r)-\Phi_N(L_N)
 =-\frac12K_Nr^2+\frac16\kappa_3r^3+R_4(N,r),
\end{equation}
where
\begin{equation}\label{app:R4-bound}
 |R_4(N,r)|\le C|K_N||r|^4.
\end{equation}
The Jacobian outside the phase contributes $e^r$.  By
\eqref{app:K-sector}, $e^{-K_Nr^2/2}$ is dominated by
$e^{-c|K_N|r^2}$.  Extending the central comparison to the full real
Gaussian introduces only an exponentially small left tail; the full real
Gaussian is not the deformed contour.  Completing the square gives
\[
 \frac{\int_{\R}r^3e^{r-K_Nr^2/2}\,\dd r}
 {\int_{\R}e^{r-K_Nr^2/2}\,\dd r}
 =\frac3{K_N^2}+\frac1{K_N^3}.
\]
Since $\kappa_3=O(|K_N|)$, the signed cubic, its square, and the quartic
remainder contribute $O(|K_N|^{-1})$ relatively.  Estimating the absolute
value of the cubic before integration would lose this order.

Outside the central window on the actual range $r\ge1-\Re L_N$, strict
horizontal concavity makes the left tail increase up to the saddle window
and gives a negative derivative of size $\gg|K_N|^{3/5}$ at the right edge.
Splitting at $|r|=1$ reduces the two resulting tail bounds to elementary
Gaussian and exponential integrals.  Together with
\eqref{app:connector-small}, this proves the leading-mode formula
\eqref{eq:leading-moment}.

\subsection{Modes and relative error}

The theta mode $m\ge2$ differs from the leading mode by
$-\pi(m^2-1)e^{L_N+r}$.  On the central window,
\eqref{app:imag-bootstrap} keeps its real part negative, while the saddle
equation gives $e^{\Re L_N}\asymp |N|/\log|N|$.  Thus
\begin{equation}\label{app:mode-suppression}
 |F_m(N)|\le|\mathcal A(N)|
 \exp\left(-c(m^2-1)|N|/\log|N|\right).
\end{equation}
The series in $m$ is dominated by its first term.  This bound is uniform on
the fixed sector; it is not asserted with a constant uniform as the sector
angle approaches $\pi/2$.

For $a(s)=\log\mathcal A(s)$, use the saddle identity during
differentiation to obtain
\[
 a'(s)=\log L_s+(L_sK_s)^{-1}-\tfrac12K_s'/K_s,
 \qquad a''(s)=O((|s|\log|s|)^{-1}).
\]
Taylor's formula on a segment of length two gives
\[
 \frac{F(N+2)}{F(N)}=L_N^2
 \left(1+O\left(\frac{\log|N|}{|N|}\right)\right).
\]
Combining this with \eqref{app:moment-assembly}, the leading-mode formula,
and sectorial Stirling proves Theorem~\ref{thm:sector}.  Holomorphy of the
relative error follows because the main term is nonzero on the outer sector
and every preceding estimate was established there before restriction to
the inner sector.

\section{Derivative and polygamma audit}\label{app:derivative-audit}

After Theorem~\ref{thm:sector}, $M_z$ is nonzero on a sufficiently remote
outer sector.  Throughout this section
\[
 h(z)=\log M_z,
\]
on the branch real on the positive axis.  In particular, $h$ is not
$\log\gamma(z)$; the gamma-ratio derivatives in the exact identity above are
kept separate.

\subsection{Implicit derivative coordinates}

The direct differentiation is organized in the two small variables
$r=L_N^{-1}$ and $\sigma=L_N/N$.  From \eqref{app:Q-factor},
\begin{equation}\label{app:Lprime-reduced}
 L_N'=\frac{1}{N(1+r-3\sigma/4)}.
\end{equation}
Every subsequent derivative of $r$ and $\sigma$ is obtained from
\eqref{app:Lprime-reduced} by the product and quotient rules.  This keeps the
denominator visible: through order six it is a power of
$4+4r-3\sigma$, whose nonvanishing has already been proved in
\eqref{app:H6-margin}.

There are two differentiation variables in the final coefficient problem.
The saddle main term is
\[
 G_0(N)=(N+1)\log L_N+\frac{L_N}{4}-\frac{N}{L_N}
 -\frac12\log Q_N.
\]
First differentiate $G_0$ with respect to $N$.  Then apply
$N=2x-2$.  The chain contributes a factor $2^j$ at derivative order $j$,
while the prefactor powers of $N$ and $L_N$ absorb the corresponding powers
in the reduced normalization.  The resulting leading constants at orders
$2,3,4,5,6$ are
\[
 2,-2,4,-12,48.
\]
At orders five and six the constants before the final chain are $-6$ and
$24$.  Recording both lists is a guard against applying the factor two only
once rather than at every differentiation.

At fifth order, setting $\sigma=0$ in the exact reduced function gives
\begin{equation}\label{app:H5-reduced}
 -\frac{6+47r+146r^2+210r^3+126r^4+42r^5+6r^6}{(1+r)^7}.
\end{equation}
At sixth order the numerator has $82$ nonzero monomials and total degree
$13$, over $(4+4r-3\sigma)^{12}$.  The canonical expanded numerator is kept
in the exact computation ledger rather than typeset twice.  A coefficientwise
triangle inequality on $|r|,|\sigma|\le7/50$ gives exactly
\begin{equation}\label{app:H6-exact-majorant}
 \frac{6422139805764931584036533551104}
 {702576099728137594188684005}<10^4.
\end{equation}
Both sides of \eqref{app:H6-exact-majorant} are rational.  No rounded decimal
is used to prove the strict inequality.

\subsection{Cauchy transport of the relative error}

Write $\gamma=\mathcal G(1+\mathcal E)$ on the outer sector.  For a positive
$x$ in the inner sector choose $R_x=\eta_0x$, with $\eta_0$ from
\eqref{app:sector-disc}.  Cauchy's formula gives
\[
 |\mathcal E^{(j)}(x)|
 \le\frac{j!}{(\eta_0x)^j}
 \sup_{|z-x|=\eta_0x}|\mathcal E(z)|.
\]
The sectorial bound on the circle is $O(\log x/x)$.  Leibniz expansion of
$\log(1+\mathcal E)$ is legal after increasing the threshold so that
$|\mathcal E|\le1/2$, and every derivative through order six is lower order
than the corresponding term of the saddle tower.  This is the step that
turns a holomorphic relative coefficient asymptotic into the uniform
logarithmic derivative estimate; differentiating a real-axis big-$O$ term
would not be valid.

\subsection{Paired polygamma differences}

For $\Re z>0$,
\[
 \psi^{(5)}(z)=120\sum_{k=0}^\infty(z+k)^{-6}.
\]
The terms are absolutely and locally uniformly convergent, so they may be
differentiated and paired on the thickened domain.  The $B-C$ contribution
is treated as one difference:
\[
 \psi^{(5)}(B+z)-\psi^{(5)}(C+z)
 =(B-C)\int_0^1\psi^{(6)}(C+z+t(B-C))\,\dd t.
\]
The $D-(n+1/2)$ pair is handled identically.  Domain containment keeps every
segment in $\Re z\ge cn$, and the series bounds the sixth polygamma by a
constant times $n^{-6}$.  The lengths of the paired segments carry the
small branch differences that would be lost by estimating the two
polygamma terms separately.  Combining these paired estimates with the
transported xi derivative proves \eqref{eq:E6}.

\section{Cube calculus and the limiting branch}\label{app:cube-branch}

\subsection{A direct proof of the cube identity}

Let $f_0(U)=\log U-\log(U+1)$.  For $q=j+1$, apply the real fundamental
theorem of calculus successively to each forward difference.  Since
\[
 \frac{\dd^q}{\dd U^q}\log U=(-1)^{q-1}(q-1)!U^{-q},
\]
Fubini gives
\[
 \Delta^jf_0(U)=(-1)^{j+1}j!
 \int_{[0,1]^q}(U+u_1+\cdots+u_q)^{-q}\,\dd u.
\]
The integrand is continuous and positive on the parameter box, so no
improper integration or interchange issue occurs.

For
\[
 \Phi_q(s,z)=\int_{[0,1]^q}
 (s+z(u_1+\cdots+u_q))^{-q}\,\dd u,
\]
differentiate under the integral.  The first two parameter derivatives are
\begin{align*}
 \partial_s\Phi_q(s,z)
 &=-q\int_{[0,1]^q}(s+z\Sigma u)^{-q-1}\,\dd u,\\
 \partial_s^2\Phi_q(s,z)
 &=q(q+1)\int_{[0,1]^q}(s+z\Sigma u)^{-q-2}\,\dd u.
\end{align*}
If $s\ge s_0>0$ and $z\ge0$, the unit cube has mass one and therefore
\begin{equation}\label{app:Phi-bounds}
 |\partial_s^r\Phi_q(s,z)|\le(q)_rs_0^{-q-r},
 \qquad r=0,1,2.
\end{equation}
Apply the mean-value theorem to the $z$ variable.  Since
$0\le\Sigma u_i\le q$,
\begin{equation}\label{app:Phi-comparison}
 \left|\partial_s^r\Phi_q(s,z)-(-1)^r(q)_rs^{-q-r}\right|
 \le(q)_r(q+r)qz\,s_0^{-q-r-1}.
\end{equation}
These two inequalities are the entire analytic content of the elementary
cube stage.

\subsection{Three limiting mechanisms}

Substitute \eqref{app:Phi-bounds}--\eqref{app:Phi-comparison} in
\eqref{eq:Hnj}.  The remote $A$ endpoint carries the external factor
$e^{q-1}$: it survives only for $q=1$.  The $B-C$ pair is a divided
difference across a segment of length $we$ and tends to $qw/t^{q+1}$.
The $D$ term must be paired with the gamma half-shift before rescaling; its
segment length tends to $\delta e-x/2$ and its normalized limit is
$q\delta$.  Keeping these mechanisms separate gives exactly
\[
 H^\infty(\alpha,t,w,\delta)=
 \left(\alpha^{-1}+wt^{-2}+\delta,
 wt^{-3}+\delta,3wt^{-4}+3\delta,4wt^{-5}+4\delta\right).
\]
The value error is $O(e+x/e)$ and the first-derivative error is $O(e+x)$,
uniformly on the rational box.  The estimates are obtained from compact
endpoint inequalities, not from a plot of the limiting branch.

\subsection{Exact limiting solution and inverse}

Adding $S^\infty=(-2,-1,-2,-2)$ gives residuals $F_0,\ldots,F_3$.  Direct
subtraction yields
\[
 3F_1-F_2=\frac{3w(t-1)}{t^4}-1,
 \qquad
 \frac43F_2-F_3=\frac{4w(t-1)}{t^5}-\frac23.
\]
At a zero, these give $3w(t-1)=t^4$ and $6w(t-1)=t^5$; positivity forces
$t=2$, then
$w=16/3$, $\delta=1/3$, and the first equation gives $\alpha=3$.  Thus the
positive solution is $y_*=(3,2,16/3,1/3)$.

The exact inverse of the Jacobian in \eqref{eq:jacobian} is
\begin{equation}\label{app:jacobian-inverse}
 P=
 \begin{pmatrix}
 -9&45&-24&9\\
 0&6&-6&3\\
 0&48&-112/3&16\\
 0&1&-4/3&1
 \end{pmatrix}.
\end{equation}
Direct multiplication in both orders gives the identity matrix, and the
largest absolute row sum is $304/3$.  For $T_n(y)=y-PG_n(y)$, the two
required certificates are
\begin{equation}\label{app:branch-certificates}
 \|G_n(y_*)\|_\infty\le\frac{3R}{608},\qquad
 \sup_{\|y-y_*\|_\infty\le R}\|DT_n(y)\|_\infty\le\frac12.
\end{equation}
The first inequality and \eqref{app:jacobian-inverse} give
$\|T_n(y_*)-y_*\|_\infty\le R/2$; the second supplies both the self-map and
the contraction.  Checking the derivative only at the center would not be
sufficient.

The outer rational box has respective left and right margins
\[
 (1/2,1/2),\quad(1/4,1/4),\quad(1/3,2/3),\quad(1/12,1/12)
\]
about $y_*$.  For $e\le1/12$, endpoint arithmetic proves the ordering
inequalities \eqref{eq:ordering-margins}.  Thus every denominator used in
the finite model is positive on the whole certified branch box.

\section{Finite-free orientation and root localization}\label{app:finitefree-details}

\subsection{The convention adapter}

Write ascending constant-term-one polynomials as
\[
 p(X)=\sum_{k=0}^d(-1)^ka_kX^k,\qquad
 q(X)=\sum_{k=0}^d(-1)^kb_kX^k.
\]
Their multiplicative finite-free convolution is
\begin{equation}\label{app:ascending-convolution}
 (p\boxconv q)(X)=
 \sum_{k=0}^d(-1)^k\binom dk^{-1}a_kb_kX^k.
\end{equation}
The cited literature states its theorems for monic descending polynomials.
Let $\operatorname{rev}_d p=X^dp(1/X)/p(0)$.  Reindexing the finite sum in
\eqref{app:ascending-convolution} gives
\[
 \operatorname{rev}_d(p\boxconv q)
 =\operatorname{rev}_d(p)\boxconv\operatorname{rev}_d(q).
\]
If $p(0)\ne0$, reversal maps every positive root to its reciprocal.  Hence
an interval $[a,b]$ is transported to $[b^{-1},a^{-1}]$, and strict
logarithmic mesh is preserved.  This proves that the orientation change
does not silently reverse either a root bound or a mesh inequality.

\subsection{Explicit MMP specialization to the comparison ${}_3F_2$}
\label{app:mmp-specialization}

This subsection spells out the specialization that is easy to obscure in a
generic ``finite-free input.''  For $U,V,s>0$ define
\begin{equation}\label{app:jacobi-factor-definition}
 Q_{U,V}^{(s)}(X)
 ={}_2F_1\!\left(\begin{matrix}-d,U\\V\end{matrix};\frac{sX}{U}\right)
 =\sum_{k=0}^d(-1)^k\binom dk
   \frac{(U)_k}{(V)_k}\left(\frac{s}{U}\right)^kX^k.
\end{equation}
Thus, in the notation of \eqref{app:ascending-convolution}, the unsigned
coefficient products of the two factors $Q_{A,B}^{(1)}$ and
$Q_{C,D}^{(D)}$ are
\[
 a_k=\binom dk\frac{(A)_k}{(B)_k}A^{-k},\qquad
 b_k=\binom dk\frac{(C)_k}{(D)_k}\left(\frac DC\right)^k.
\]
Substitution in \eqref{app:ascending-convolution} gives
\begin{align}
 [Q_{A,B}^{(1)}\boxconv Q_{C,D}^{(D)}](X)
 &=\sum_{k=0}^d(-1)^k\binom dk
   \frac{(A)_k(C)_k}{(B)_k(D)_k}
   \left(\frac{D}{AC}\right)^kX^k\nonumber\\
 &={}_3F_2\!\left(
   \begin{matrix}-d,A,C\\B,D\end{matrix};\frac{D}{AC}X\right)
 =p_F(X).                                      \label{app:concrete-3f2-convolution}
\end{align}
This is an identity of finite coefficient lists, with no root theorem used.
The Lean theorem
\code{terminating3F2Polynomial\_eq\_finiteFree\_jacobiFactors} proves the
same identity for symbolic positive $A,B,C,D$, and
\code{xiNaturalComparisonPolynomial\_eq\_finiteFree} substitutes the actual
xi branch parameters.

It remains to check that the two factors in
\eqref{app:concrete-3f2-convolution} are in MMP's positive-rooted class.  The
standard Jacobi identity is
\begin{equation}\label{app:jacobi-identification}
 Q_{U,V}^{(s)}(X)
 =\frac{d!}{(V)_d}
 P_d^{(V-1,\,U-V-d)}\!\left(1-\frac{2sX}{U}\right).
\end{equation}
On the branch used in the proof,
\[
 B,D>0,\qquad A\ge B+d,\qquad C\ge D+d,
 \qquad 1,D>0.
\]
For either factor in \eqref{app:jacobi-identification}, the Jacobi parameters
$V-1$ and $U-V-d$ are therefore greater than $-1$.  The classical Jacobi
zero theorem gives $d$ simple zeros in $(-1,1)$.  The affine map in
\eqref{app:jacobi-identification} sends them to
\[
 X_j=\frac{U}{2s}(1-x_j)>0.
\]
The top coefficient in \eqref{app:jacobi-factor-definition} is nonzero, so
each factor has degree exactly $d$; its constant term is one.  These facts
verify membership in the positive-rooted degree-$d$ class required by MMP,
not merely numerical plausibility for selected parameters.

Now reflect both factors to MMP's monic descending convention.  The reflection
identity proved above shows that their MMP convolution is precisely the
reflection of the left side of
\eqref{app:concrete-3f2-convolution}.  MMP v3 Proposition~2.7(iii) places all
convolution roots in $[0,\infty)$.  Since the ascending constant term is one,
zero is excluded.  For $d\ge2$, the simple positive roots of the first Jacobi
factor have
\[
 \operatorname{lmesh}(Q_{A,B}^{(1)})>1.
\]
MMP v3 Definition~2.16 and Proposition~2.17 give
\[
 \operatorname{lmesh}(Q_{A,B}^{(1)}\boxconv Q_{C,D}^{(D)})
 \ge \operatorname{lmesh}(Q_{A,B}^{(1)})>1,
\]
so the product roots are distinct.  For $d=0$ the assertion is vacuous, and
for $d=1$ the nonzero linear coefficient gives one simple positive root.
Equation \eqref{app:concrete-3f2-convolution} finally transfers the conclusion
to the paper's particular ${}_3F_2$.

{\setlength{\emergencystretch}{1em}
The formal boundary now mirrors exactly this calculation.
{\small\code{MMPFiniteFreeLogMeshInput}} takes the two displayed factor polynomials,
their complete positive-root and degree certificates, and the MMP conclusion
for their finite-free convolution.  It does not take the final xi comparison
function as an assumed positive-rooted object.
Lean then rewrites the convolution to that function using the independently
proved coefficient identity.  The general Jacobi zero theorem and the two
general MMP propositions remain declared literature inputs; the entire
specialization from those inputs to this ${}_3F_2$ is explicit.
\par}

\subsection{A ratio-free Jacobi interval}

Transport the classical Jacobi matrix from $[-1,1]$ to the positive-root
coordinate.  Under $U\ge V+d$ and $V\ge32d$, its diagonal displacement and
adjacent off-diagonal row sum satisfy
\[
 |J_{ii}-V|\le4\sqrt{Vd},\qquad
 |J_{i,i-1}|+|J_{i,i+1}|\le4\sqrt{Vd}.
\]
Gershgorin's theorem therefore places every root in
\begin{equation}\label{app:jacobi-eight}
 [V-8\sqrt{Vd},\,V+8\sqrt{Vd}].
\end{equation}
This estimate does not assume that the ratio of the two finite model gaps is
small.  That distinction matters because the new branch has
$(C-D)/D\to1$, whereas the sharper estimate used in the earlier argument
requires this ratio to be at most $1/4$.

The preliminary inequalities in \eqref{eq:coarse-parameter} make both
Jacobi intervals positive.  In particular the safe threshold is
$K_{\rm pre}=256$: using $32$ in the generic second factor would give the
lower relative endpoint $1-8/\sqrt{32}<0$.  Applying the maximum-root
product theorem in the original and reciprocal orientations gives both
endpoints of the convolution interval.  If the deviations of the two
factors are $u$ and $v$, then
\begin{align*}
 |(B+u)(1+v)-B|
 &\le8\sqrt{Bd}+8B\sqrt{d/D}
       +64\sqrt{Bd}\sqrt{d/D}\\
 &\le(12+8\sqrt6)\sqrt{Bd}.
\end{align*}
This is $C_{\rm loc}=12+8\sqrt6$, not the transposed expression
$8+12\sqrt6$.

\section{Recurrence and global-maximum argument}\label{app:recurrence-details}

\subsection{Finite coefficient proof}

For fixed $m$, initialize $c_0=1$ and impose the exact finite recurrence
\begin{equation}\label{app:finite-coefficients}
 AC(k+1)(k+B+m)(k+D+m)c_{k+1}
 =D(k+m-d)(k+A+m)(k+C+m)c_k.
\end{equation}
At $k=d-m$ the right side vanishes, so the producer terminates.  Every
denominator on the preceding range is positive.  Equation
\eqref{app:finite-coefficients} therefore proves the coefficient ratio
without asking a symbolic system to simplify a quotient of analytically
continued Pochhammer functions at unspecified arguments.

On a monomial $y^k$, the Euler operator $\vartheta=y\,\dd/\dd y$ acts by
multiplication by $k$.  Applying both sides of \eqref{eq:euler-ode} to the
finite coefficient list and using \eqref{app:finite-coefficients} proves the
ODE coefficient by coefficient.  The shift identity
$(a)_{k+m}=(a)_m(a+m)_k$ identifies the shifted polynomial with the $m$th
derivative up to a nonzero factor.

Expanding the ODE gives the equivalent coefficients
\begin{align}
 P_{3,m}&=\frac{AC-Dy}{AC},\nonumber\\
 P_{2,m}&=B+D+2m+1
 -\frac{Dy}{AC}(A+C+3m-d+3),\nonumber\\
 P_{1,m}&=(B+m)(D+m)
 -\frac{Dy}{AC}\bigl((A+m)(C+m)
 +(m-d)(A+C+2m+1)\bigr),\label{app:P-expanded}\\
 P_{0,m}&=-\frac{Dy}{AC}(m-d)(A+m)(C+m).\nonumber
\end{align}
The final coefficient is nonnegative for $0\le m\le d$ and positive before
termination.  Exact algebra proves equality of \eqref{app:P-expanded} with
the decomposed formulas \eqref{eq:Pcoeffs}; the decomposed form, rather than
the fully expanded form, is used for inequalities.

\subsection{Normalization of the radius}

At a critical point set
\[
 T_k=\frac{y^kp_F^{(k)}(y)}{p_F(y)}.
\]
Then $T_0=1$ and $T_1=0$.  There is no $1/k!$ in this definition.  Divide
the recurrence at level $m$ by $P_{2,m}R^{m+2}$, where
$R=K_r\sqrt{Bd}$ and $K_r=4096$.  Under the exact bounds
\eqref{eq:recurrence-bounds}, the contributions of $P_0$, $P_1$, and $P_3$
are bounded by fixed strict fractions of one.  The remaining inequalities
use only the certified caps on $d/n$, $B/n$, and $D/n$.

Set $\mathcal M=\max_{0\le j\le d}|T_j|/R^j$.  If $\mathcal M>1$, choose a
maximizing index $k$.  Since $T_0=1$ and $T_1=0$, we have $k\ge2$.  Use the
recurrence with $m=k-2$, isolate its $P_{2,k-2}T_k$ term, and divide by
$P_{2,k-2}R^k\mathcal M$.  Every other normalized neighbor is at most one by
the definition of $\mathcal M$; if $k=d$, the higher neighbor is
$T_{d+1}=0$ by termination.  The three coefficient budgets have sum strictly
below one, contradicting maximality.  Therefore
\[
 \max_{1\le k\le d}|T_k|^{1/k}\le K_r\sqrt{Bd}.
\]
This global maximum, rather than lower-index minimality, is what controls the
higher neighbor $T_{k+1}$.  Thus no boundary order is silently omitted.

\section{Six-node interpolation and stability}\label{app:HG-details}

\subsection{The complex line-segment recursion}

For holomorphic $f$ on an open convex set and $a,b$ in that set, complex FTC
along the segment gives
\begin{equation}\label{app:complex-ftc}
 f(b)-f(a)=(b-a)\int_0^1f'(a+t(b-a))\,\dd t.
\end{equation}
Apply \eqref{app:complex-ftc} recursively to Newton divided differences.
After six steps, every evaluation point is a convex combination of
$z_0,\ldots,z_5,z$.  A stick-breaking parametrization of the standard
six-simplex produces the weight
\[
 (1-u_0)^5(1-u_1)^4(1-u_2)^3(1-u_3)^2(1-u_4).
\]
Its integral over $[0,1]^6$ is
\[
 \frac16\frac15\frac14\frac13\frac12=\frac1{720}.
\]
If the six node values vanish, the Newton formula therefore reduces to
\[
 f(z)=\prod_{j=0}^5(z-z_j)
 \int_{\Delta_6}f^{(6)}(\zeta(t))\,\dd t,
\]
which proves Lemma~\ref{lem:HG}.  Convexity is used exactly once: it keeps
$\zeta(t)$ inside the derivative-bound domain.

\subsection{Why the sixth match gives the exponent}

For $E_F$, the nodes are $0,1,\ldots,5$.  On the thickened critical-point
domain, projection to $[0,d]$ and $\rho\ge d$ give $|z-j|\le3\rho$.
Consequently
\begin{align*}
 |E_F(z)|
 &\le\frac{C_{B6}(3\rho)^6}{720n^5\log(n+2)}\\
 &\le C_\ast\frac{d^3}{n^2\log(n+2)},
\end{align*}
because $\rho=K_r\sqrt{Bd}$ and $B\le3n$.  Five matched nodes would produce
a fifth power of the localization radius and recover the older three-fifths
scale.  The gain to two-thirds is therefore exactly the extra matching
condition together with one additional derivative estimate.

The logarithm branch is legitimate because the coefficient ratios are
nonzero on the parameter and complex-domain boxes.  Once $|E_F|\le1$, the
segment identity
\[
 e^{E_F}-1=E_F\int_0^1e^{tE_F}\,\dd t
\]
gives $|e^{E_F}-1|\le e|E_F|$; the final constant is enlarged to use the
simpler bound $2|E_F|$ on the smaller certified disc.  This converts the
logarithmic interpolation estimate to the multiplier estimate required in
Lemma~\ref{lem:stability}.

At each critical point, Cauchy's formula on the circle of radius $\rho$
turns the bound on $c_F-1$ into a convergent derivative tail.  The critical
point radius of Proposition~\ref{prop:radius} supplies the matching bound on
the model derivatives.  Rouch\'e on disjoint critical-point discs then
preserves one simple zero in each gap.  Positivity of coefficients and the
constant-term convention fix the sign of those roots; the rescaling
$y=-SX$ sends them to distinct negative roots of $\J^{d,n}$.

\section{Effectivity, dependencies, and logical boundary}\label{app:effectivity-details}

\subsection{The finite pre-cutoff range is empty}

Let $N_0$ be an integer above both the analytic and explicit construction
cutoffs and set
\[
 K_{\rm finite}=N_0^2(N_0+2)+1.
\]
For natural $0\le n<N_0$, the elementary inequality
$\log x\le x-1$ gives
\[
 n^2\log(n+2)\le n^2(n+1)<N_0^2(N_0+2)<K_{\rm finite}.
\]
If $d\ge1$, then $K_{\rm finite}\le K_{\rm finite}d^3$.  Hence
\[
 K_{\rm finite}d^3>n^2\log(n+2),
\]
which is the negation of the wedge antecedent.  Increasing the final constant
above $K_{\rm finite}$ only strengthens this exclusion.  Thus ``finite
absorption'' here does not mean that finitely many unverified Jensen
polynomials are being declared hyperbolic; it means that the stated
implication has no antecedent below the analytic cutoff.

The Lean definition \code{finiteCutoffAbsorptionConstant} and the theorem
whose source name is \code{not\_twoThirdsWedge\_} followed by
\code{finiteCutoffAbsorption} formalize this calculation over natural
$N_0,n,d$.  This closes the formerly paper-only pre-cutoff step.

\subsection{Finite and analytic constants}

Let $x_{\rm admissible}$ be the smallest rational cap on $d/n$ demanded by
the branch box, the Jacobi intervals, the global-maximum budgets, and the
complex-domain containment.  Since $\log(n+2)\le n$ for $n\ge2$, the wedge
implies $d/n\le K^{-1/3}$.  The exact finite ledger takes
\begin{equation}\label{app:K-geometry}
 K_{\rm geometry}
 =\lceil x_{\rm admissible}^{-3}\rceil+1
 =522255954073519803473068032000000000000000001.
\end{equation}
The size of \eqref{app:K-geometry} reflects deliberately coarse uniform
inequalities; it is not a prediction of the true transition.

Only two analytic quantities remain symbolic:
\begin{enumerate}
\item $N_{\rm analytic}$, the threshold after which the fixed-sector saddle,
contour, and coefficient estimates all hold;
\item $C_{B6}$, the explicit constant in the uniform sixth-derivative bound
\eqref{eq:E6}.
\end{enumerate}
Once those are read from the displayed paper estimates, the final proof may
take the maximum of $K_{\rm geometry}$, the threshold conversion, and the
multiplier constant.  Thus the theorem is effective even though this paper
does not optimize or print a practically useful numerical $K$.

\subsection{What is and is not kernel checked}

The Lean development proves the concrete centered-xi theta/Mellin/omega
identity, factor eight, quantitative moving-sector saddle, legal contour,
Gaussian localization, infinite higher-mode suppression, Gamma/Stirling
assembly, and uniform coefficient theorem through order six.  It also proves
the exact xi whole-box branch, transformed comparison, six coefficient
matches, finite-free specialization, complete critical-point localization,
and literal normalized radius estimate.  The Phase-26 axiom audit covers
1,217 declarations and reports only standard logical principles:
propositional extensionality, classical choice, and quotient soundness.  The
project contains no \code{sorry}, \code{admit}, project-specific axiom, or
unsafe escape hatch on this proof surface.

The Phase-28 terminal audit additionally covers the effective and literal
forms of Theorem~\ref{thm:sector} and its derivative consequence through
order six.  The Phase-29 Palomar package isolates these statements and has
passed the repository's local fresh-kernel, source-fidelity, and mutation
checks.  Palomar's own Comparator and NanoDa services have not yet replayed
the candidate, so this document makes no official Palomar-verification
claim.

The Phase-30 endpoint constructs the final xi-specific sixth-multiplier
interval certificate, instantiates the sign-transfer theorem, identifies the
actual transformed xi Jensen polynomial, and proves the headline
negative-root conclusion.  Exact symbolic systems and Arb/ACB remain
corroborating channels rather than proof premises.  The mathematical
statements imported from classical Jacobi theory, MMP, and MSS are typed
inputs to the concrete headline theorem rather than re-proved in Lean.

The most heavily imported external results are the classical Jacobi
matrix/root correspondence, the finite-free nonnegative-root and log-mesh
theorems in the pinned MMP version, and the MSS maximum-root theorem.  The
source audit records exact versions and theorem numbers, translates every
polynomial convention explicitly, and marks these results as imported rather
than re-proved.  This separation is essential: a verified adapter is not a
formalization of the external theorem it adapts.

\subsection{Review status}

The available analytic and algebraic adversarial reports were produced by
AI systems with separated prompts and independent-recalculation requirements.
They have found material transcription errors during development, and those
errors were corrected and protected by regression tests.  They are useful
evidence, but they are neither human nor peer review.  The release materials
state this limitation prominently so that a prospective reviewer can assign
the appropriate evidentiary weight to every channel.
